\itemhead{复习题三 · 第4题（教材p.90）}{XJ-B1-C03-RV-P090-Y0720}
\label{XJ-B1-C03-RV-P090-Y0720}
{\footnotesize 湘教版教材p.90 / PDF第97页 / 复习题三 / 第4题}
\field{原题}{已知函数 $y=f(x)$ 的对应关系如下表，函数 $y=g(x)$ 的图象是如图所示的曲线 $ABC$，其中 $A(1,3),B(2,1),C(3,2)$，则 $f(g(2))$ 的值为\blank。
\begin{center}\begin{tabular}{c|rrr}\hline$x$&1&2&3\\$f(x)$&2&3&0\\\hline\end{tabular}\end{center}
\bankfig{C03-RV090}{0.38}}
\field{参考答案}{$g(2)=1,f(1)=2$，故答案为2。}
{\small \field{关联知识}{F03 函数求值、解析式与函数建模（教材帮p.80、87、88、89、90、131 / PDF89、96、97、98、99、140）；F04 函数三种表示法、图象与读数（教材帮p.85、86、93、94 / PDF94、95、102、103）}
\field{方法与易错点}{标清每次输入和输出；按已知函数类型或代换关系建立等价方程。 易错：将 $f(x)$ 当成 $f	imes x$；复合顺序反；换元后沿用错误定义域。}
\field{答案依据}{答案册相关题组：PDF第24、25、26、27页（册内页23、24、25、26）；按当前题库原题逐项重算，不按旧版题号机械对套。}}
